% vim:ts=4:sw=4
%
% Copyright (c) 2008-2009 solvethis
% Copyright (c) 2010-2012 Casper Ti. Vector
% Public domain.
%
% 使用前请先仔细阅读 pkuthss 和 biblatex-caspervector 的文档，
% 特别是其中的 FAQ 部分和用红色强调的部分。
% 两者可在终端/命令提示符中用
%   texdoc pkuthss
%   texdoc biblatex-caspervector
% 调出。

% 在黑白打印时彩色链接可能变成浅灰色，此时可将“colorlinks”改为“nocolorlinks”。
\documentclass[adobefonts,UTF8, colorlinks, notocbibind]{pkuthss}

% 使用 biblatex 排版参考文献，并规定其格式。
\usepackage[backend = biber, style = caspervector, utf8]{biblatex}
% 使得打字机粗体可以被使用。
\usepackage{lmodern}
% 产生 originauth.tex 里的 \square。
\usepackage{amssymb}
% 提供 Verbatim 环境和 \VerbatimInput 命令。
\usepackage{fancyvrb}
% 图像包
\usepackage{subcaption}
% --引入其他pdf文件--
\usepackage{pdfpages}

% 使被强调的内容为红色。
\newcommand{\myemph}[1]{\emph{\textcolor{red}{#1}}}

% pkuthss 文档模版的版本。
\newcommand{\docversion}{v1.4 rc1}

% 设定文档的基本信息。
\pkuthssinfo{
	cthesisname = {本科生毕业论文}, ethesisname = {Undergraduate Thesis},
	ctitle = {中山大学论文文档模版\\pkuthss \docversion},
	% “\\”在设定 pdf 属性时会被自动过滤掉，于是得到的 pdf 属性中标题为
	%   The PKU dissertation document classpkuthss [版本号]
	% 此处指定其被替换为“: ”，以使之为
	%   The PKU dissertation document class: pkuthss [版本号]
	etitle = {%
		The PKU dissertation document class\texorpdfstring{\\}{: }%
		pkuthss \docversion%
	},
	cauthor = {罗尼 $\cdot$ 奥沙利文},
	eauthor = {Ronnie O'Sullivan},
	studentid = {07302???},
	date = {二〇一一年五月},
	school = {信息科学与技术学院},
	cmajor = {计算机科学与技术}, emajor = {Computer Science and Technology},
	direction = {强力斯诺克},
	cmentor = {XX 教授}, ementor = {Prof.\ XX},
	ckeywords = {\LaTeX2e{}，排版，文档类，\CTeX{}},
	ekeywords = {\LaTeX2e{}, typesetting, document class, \CTeX{}}
}
% 导入参考文献数据库（注意不要省略“.bib”）。
\addbibresource{sample.bib}

\begin{document}
	% 以下为正文之前的部分。
	\frontmatter

	% 自动生成标题页。
	\maketitle
	% --可以使用includepdf直接将转换好的附表或封面加入--
	%\includepdf[pages=1-2]{chap/cover.pdf}
	%\includepdf[pages=1-2]{chap/report.pdf}
	%\includepdf[pages=1-2]{chap/check.pdf}
	%\includepdf[pages=1-2]{chap/reply.pdf}
	%\includepdf[pages=1-2]{chap/statement.pdf}
	% 版权声明。
	% vim:ts=4:sw=4
%
% Copyright (c) 2008-2009 solvethis
% Copyright (c) 2010-2012 Casper Ti. Vector
% All rights reserved.
%
% Redistribution and use in source and binary forms, with or without
% modification, are permitted provided that the following conditions are
% met:
%
% * Redistributions of source code must retain the above copyright notice,
%   this list of conditions and the following disclaimer.
% * Redistributions in binary form must reproduce the above copyright
%   notice, this list of conditions and the following disclaimer in the
%   documentation and/or other materials provided with the distribution.
% * Neither the name of Peking University nor the names of its contributors
%   may be used to endorse or promote products derived from this software
%   without specific prior written permission.
% 
% THIS SOFTWARE IS PROVIDED BY THE COPYRIGHT HOLDERS AND CONTRIBUTORS "AS
% IS" AND ANY EXPRESS OR IMPLIED WARRANTIES, INCLUDING, BUT NOT LIMITED TO,
% THE IMPLIED WARRANTIES OF MERCHANTABILITY AND FITNESS FOR A PARTICULAR
% PURPOSE ARE DISCLAIMED. IN NO EVENT SHALL THE COPYRIGHT HOLDER OR
% CONTRIBUTORS BE LIABLE FOR ANY DIRECT, INDIRECT, INCIDENTAL, SPECIAL,
% EXEMPLARY, OR CONSEQUENTIAL DAMAGES (INCLUDING, BUT NOT LIMITED TO,
% PROCUREMENT OF SUBSTITUTE GOODS OR SERVICES; LOSS OF USE, DATA, OR
% PROFITS; OR BUSINESS INTERRUPTION) HOWEVER CAUSED AND ON ANY THEORY OF
% LIABILITY, WHETHER IN CONTRACT, STRICT LIABILITY, OR TORT (INCLUDING
% NEGLIGENCE OR OTHERWISE) ARISING IN ANY WAY OUT OF THE USE OF THIS
% SOFTWARE, EVEN IF ADVISED OF THE POSSIBILITY OF SUCH DAMAGE.

\chapter*{版权声明}
{
	\zihao{3}\linespread{1.5}\selectfont

\iffalse
	任何收存和保管本论文各种版本的单位和个人，
	未经本论文作者同意，不得将本论文转借他人，
	亦不得随意复制、抄录、拍照或以任何方式传播。
	否则一旦引起有碍作者著作权之问题，将可能承担法律责任。
	\par
\fi
	版权所有 \copyright\ 2008--2009 solvethis
	\par
	版权所有 \copyright\ 2010--2012 Casper Ti. Vector
	\vskip 1em

	pkuthss 文档类和 pkuthss-extra 宏包以 %
	\LaTeX{} Project Public License 发布。
	本说明（示例）文档的源代码中，除以下文件
	\begin{itemize}
		\item \verb|img/sysulogo.eps|
		\item \verb|chap/copyright.tex|
		\item \verb|chap/originauth.tex|
	\end{itemize}
	以 New BSD License 发布，
	以及参考文献和引用使用的样式文件版权声明见相应文件中说明以外，
	其余部分文件发布在公有领域（public domain）。
	\par
}


	% 中英文摘要。
	% vim:ts=4:sw=4
%
% Copyright (c) 2008-2009 solvethis
% Copyright (c) 2010-2012 Casper Ti. Vector
% Public domain.

\begin{cabstract}

	本文介绍了 pkuthss 文档模版所提供的功能，
	并以自身为例演示了该模版的使用。

\end{cabstract}

\begin{eabstract}

	This paper describes the the functions provided by
	the pkuthss document template,
	and provides itself as an example to illustrate
	the usage of the document class.

\end{eabstract}


	% 自动生成目录。
	\tableofcontents

	% 以下为正文。
	\mainmatter

	% 绪言。
	% vim:ts=4:sw=4
%
% Copyright (c) 2008-2009 solvethis
% Copyright (c) 2010-2012 Casper Ti. Vector
% Public domain.

\specialchap{绪言}

本文档是“北京大学论文文档模版”的说明文档，
同时也是使用模版的一个示例。

pkuthss 文档模版由三部分构成：
\begin{itemize}
	\item \textbf{pkuthss 文档类}：
		其中进行了学位论文所需要的一些基本的设定，
		主要包括对基本排版格式的设定和提供设置论文信息的命令。
	\item \textbf{pkuthss-extra 宏包}：
		其中实现了学位论文中用户可能较多用到的一些额外功能，
		例如自动在目录中加入参考文献和索引的条目和%
		自动根据用户设定的文档信息对所生成 pdf 的作者、标题等属性进行设置等。
	\item \textbf{说明（示例）文档}：
		说明文档即本文档，
		在安装（见第 \ref{sec:inst} 节）之后应该可以用 \TeX{} 系统提供的
		\verb|texdoc| 命令调出：
\begin{Verbatim}[frame=single]
texdoc pkuthss
\end{Verbatim}
		同时，
		本文档的源代码（位和本文档的 pdf 文件处于同一目录下）%
		也正是用户撰写自己的学位论文时的一个模版：
		用户只需按照模版中的框架修改代码，
		即可写出自己的论文。
\end{itemize}

在此之前，包括 dypang\supercite{dypang}、FerretL\supercite{FerretL}、%
lwolf\supercite{lwolf}、Langpku\supercite{Langpku}、%
solvethis\supercite{solvethis} 等的数位网友均做过学位论文模版的工作。
本论文模版是 solvethis 的 pkuthss 模版的更新版本，
更新的重点是重构和对新文档类、宏包的支持。

pkuthss 文档模版现在的维护者是 Casper Ti. Vector\footnote%
{\href{mailto:CasperVector@gmail.com}{\texttt{CasperVector@gmail.com}}}。%
pkuthss 文档模版目前托管在 Google Code 上，
其项目主页是：\\
\hspace*{\parindent}\url{http://code.google.com/p/caspervector/}


	% 各章节。
	% vim:ts=4:sw=4
%
% Copyright (c) 2008-2009 solvethis
% Copyright (c) 2010-2012 Casper Ti. Vector
% Public domain.

\chapter{使用介绍}
	\section{系统要求}\label{sec:req}

	正确编译需要以下几部分：
	\begin{itemize}
		\item 一个基本的 \TeX{} 发行版。
		\item CJK 或 XeCJK（供 Xe\LaTeX{} 使用）宏包。
		\item ctex 宏包\supercite{ctex,ctex-faq}（提供了 ctexbook 文档类）。
		\item 中文字体。
		\item 如果要使用 biblatex 进行文献列表和引用的排版的话，
			还需要 biblatex 宏包\supercite{biblatex}；
		\item 如果需要对中文文献进行按汉语拼音的排序的话，
			还需要 biber 程序\supercite{biber}。
		\item 如果默认的文献列表和引用样式的话，
			还需要作者编写的 biblatex 样式
			（biblatex-caspervector）\supercite{biblatex-caspervector}。
		\item 如果需要使用 Makefile 来实现自动编译，还需要 Make 工具；
			但如果使用由批处理实现的伪“Makefile”就不用了。
	\end{itemize}

	\myemph{最新}的\myemph{完全版} \TeX{}Live 系统（\myemph{%
		注：
		某些 Linux 发行版软件仓库中的 \TeX{}Live 有问题，
		建议使用独立安装版的 \TeX{}Live%
		\footnote{\url{http://www.tug.org/texlive/}.}。%
	}）和 \CTeX{} 套装都已经包含除中文字体和 Make 之外所有要求的项目。%
	\myemph{%
		为了获得最好的支持，
		我们建议用户使用最新完全版的 \TeX{} 系统和各宏包。%
	}根据目前已有的测试结果，
	使用 \TeX{}Live 2009 及其以上稳定版本%
	或 \CTeX{} 套装 v2.8 系列及其以上版本
	的用户可以正常使用本模版。

	中文字体需要用户自行获得。\myemph{%
		注：
		一些中文字体的字库不全，
		只有 GB2312 字符集内字符的字体信息。
		这种情况通常会造成编译生成的 pdf 文件中缺少部分字符，
		其中一种典型症状是封面的“〇”字显示不出来。
		如果要使用中易公司的字体，
		则建议使用 Windows Vista 及其以后版本提供的%
		宋体、黑体、楷体和仿宋体，
		以及 Microsoft Office 2003 及其以后版本提供的隶书和幼圆体，
		这些字体是 GB18030 字符集的，不存在上述问题。%
	}

	Linux 用户可以从软件源获得 GNU 的 make；
	其它类 UNIX 系统应该也会提供 make 工具，请参阅相应的文档以获得帮助。%
	Windows 用户可以从以下地址下载 Windows 下的 GNU make 工具：\\
	\hspace*{\parindent}%
	\url{http://gnuwin32.sourceforge.net/packages/make.htm}

	\section{安装方法}\label{sec:inst}

	使用 \TeX{}Live 的用户可以通过在终端中以管理员权限执行
\begin{Verbatim}[frame = single]
tlmgr install pkuthss
\end{Verbatim}
	来安装 pkuthss 文档模版。

	使用 \CTeX{} 套装的用户可以通过在命令提示符中执行
\begin{Verbatim}[frame = single]
mpm --install=pkuthss
\end{Verbatim}
	来安装 pkuthss 文档模版。

	\section{模版文件}\label{sec:doc-dir}

	在正确安装 pkuthss 文档模版之后，在终端/命令提示符中执行
\begin{Verbatim}[frame = single]
texdoc pkuthss
\end{Verbatim}
	所打开的 pdf 所在的同一目录中包含本文档（\verb|pkuthss.pdf|）的源代码%
	（\verb|utf8lf/| 和 \verb|gbkcrlf/| 两个子目录，
	两目录中代码除使用的编码和换行符外基本上完全相同）。

	其中，%
	\verb|utf8lf/| 目录下的源代码使用的是 UTF-8 编码、\verb|\n|（LF）换行，
	适合类 UNIX 系统用户使用；%
	\verb|gbkcrlf/| 目录下的源代码使用的是 GBK 编码、\verb|\r\n|（CRLF）换行，
	适合 Windows 用户使用。%
	用户可以试情况将 \verb|utf8lf/| 或 \verb|gbkcrlf/| 中
	的所有内容复制到合适的目录，
	并在此目录中根据模版修改出自己的论文。

	\verb|utf8lf/| 和 \verb|gbkcrlf/| 目录中的重要文件有：
	\begin{itemize}
		\item \verb|Makefile|：
			被 Make 工具调用的 Makefile，用于使编译工作自动化。
		\item \verb|Make.bat|：%
			Windows 下的伪“Makefile”，用 Windows 批处理实现。
		\item \verb|chap/|：目录，包含各章节内容：
		\begin{itemize}
			\item \verb|copyright.tex|：
				版权声明部分\footnote{%
					本文档中的版权声明并不是%
					学校默认要求的形式\supercite{pku-copyright}。
					符合学校要求的一个版权声明已经放在此文件中，
					但用 \texttt{\string\iffalse{} ...\ \string\fi} %
					注释掉了，
					用户可以考虑使用那个版本。
				}。
			\item \verb|originauth.tex|：
				原创性声明和使用授权说明部分\supercite{pku-originauth}。
		\end{itemize}
		\item \verb|img/|：目录，包含论文中所有图片：
		\begin{itemize}
			\item \verb|Makefile|：图片部分的 Makefile。
			\item \verb|Make.bat|：%
				Windows 下的伪“Makefile”，用 Windows 批处理实现。
			\item \verb|pkulogo.eps|：北京大学校徽。
			\item \verb|pkuword.eps|：“北京大学”字样。
		\end{itemize}
			\myemph{%
				注：%
				pdf\LaTeX{} 方式编译（见第 \ref{sec:compile} 节）
				可能需要将 eps 图片转换为 pdf 格式，
				而使用 Makefile 或伪“Makefile”时这些图片可以自动生成；
				不使用 Makefile 的用户可以手动运行 %
				\texttt{img/} 目录中的 Makefile 或伪“Makefile”
				来生成这些图片。%
			}
	\end{itemize}

	\section{编译方式}\label{sec:compile}

	pkuthss 文档模版支持三种编译方式，即
	\begin{itemize}
		\item \LaTeX{} -- dvipdfmx 方式：\\
			依次执行 \verb|latex|，\verb|biber|（或 \verb|bibtex|），%
			\verb|latex|，\verb|latex| 和 \verb|dvipdfmx|。
		\item pdf\LaTeX{} 方式：\\
			依次执行 \verb|pdflatex|，\verb|biber|（或 \verb|bibtex|），%
			\verb|pdflatex| 和 \verb|pdflatex|。
		\item Xe\LaTeX{} 方式：\\
			依次执行 \verb|xelatex|，\verb|biber|（或 \verb|bibtex|），%
			\verb|xelatex| 和 \verb|xelatex|。\\
			\myemph{%
				注意：Xe\LaTeX{} 对非 UTF-8 的编码支持不好，
				因此 Xe\LaTeX{} 方式的编译不支持 GBK 编码。
			}
	\end{itemize}

	pkuthss 文档模版附带的 Makefile 中已经对这三种编译方式进行了完整的配置。
	用户只需要在 Makefile 中通过设定变量 \verb|JOBNAME| 的值%
	指定被编译的主文件名，
	并通过设定变量 \verb|LATEX| 的值指定采用哪种编译方式，
	即可通过在主文件所在目录调用 Make 工具来实现自动编译：
	如果是在类 UNIX 环境下，则用户应该调用的命令名为 \verb|make|；
	而如果是在 Windows 环境下，
	则用户应该调用的命令名可能为 \verb|mingw32-make|。

	用户如果不想配置 Windows 下的 GNU Make，
	则也可以使用由 Windows 批处理实现的伪“Makefile”，
	通过在主文件所在目录调用 \verb|make|\footnote{%
		Windows 将批处理文件作为可执行文件，
		调用时可以不显式地指出扩展名。%
	} 或直接双击 \verb|Make.bat| 的图标运行之。


	% vim:ts=4:sw=4
%
% Copyright (c) 2008-2009 solvethis
% Copyright (c) 2010-2012 Casper Ti. Vector
% Public domain.

\chapter{pkuthss 文档模版提供的功能}
	\section{pkuthss 文档模版提供的文档类和宏包选项}
		\subsection{pkuthss 文档类提供的选项}\label{ssec:options}

		\begin{itemize}
			\item \textbf{\texttt{[no]extra}}：
				用于确定是否自动载入 pkuthss-extra 宏包。
				在默认情况下，pkuthss 文档类将使用 \verb|extra| 选项。
				用户如果不需要自动载入 pkuthss-extra 宏包，
				则需要在载入 pkuthss 时加上 \verb|noextra| 选项。

			\item \textbf{pkuthss-extra 宏包提供的选项}：
				这些选项将被传递给 pkuthss-extra 宏包
				（用户需要启用 \verb|extra| 选项）。
				具体说明参见第 \ref{ssec:extra} 小节。

			\item \textbf{其余文档类选项}：%
				pkuthss 文档类以 ctexbook 文档类为基础，
				其接受的其余所有文档类选项均被传递给 ctexbook。
				其中可能最常用的选项是 \verb|GBK| 和 \verb|UTF8|：
				它们选择源代码使用的字符编码，默认使用 \verb|GBK|。
		\end{itemize}

		例如，如果需要使用 UTF-8 编码撰写论文，
		则需要在导入 pkuthss 文档类时加上 \verb|UTF8| 选项：
\begin{Verbatim}[frame = single]
\documentclass[UTF8, ...]{pkuthss} % “...”代表其它的选项。
\end{Verbatim}

		又例如，文档默认情况下是双面模式，每章都从右页（奇数页）开始。
		如果希望改成一章可以从任意页开始，可以这样设置：
\begin{Verbatim}[frame = single]
\documentclass[openany, ...]{pkuthss} % 每章从任意页开始。
\end{Verbatim}
		但这样设置时左右（奇偶）页的页眉页脚设置仍然是不同的。
		如果需要使左右页的页眉页脚设置一致，可以直接采用单面模式：
\begin{Verbatim}[frame = single]
% 使用 oneside 选项时不需要再指定 openany 选项。
\documentclass[oneside, ...]{pkuthss}
\end{Verbatim}

		\subsection{pkuthss-extra 宏包提供的选项}\label{ssec:extra}

		除非特别说明，
		下面提到的选项中都是不带“\verb|no|”的版本被启用。

		\begin{itemize}
			\item \textbf{\texttt{[no]spacing}}：
				是否采用一些常用的对空白进行调整的版式设定。
				具体地说，启用 \verb|space| 选项后会进行以下几项设置：
			\begin{itemize}
				\item 自动忽略 CJK 文字之间的空白而%
					保留（CJK 文字与英文之间等的）其它空白。
				\item 调用 setspace 宏包以使某些细节处的空间安排更美观。
				\item 设置页芯居中。
				\item 设定行距为 1.41\footnote{%
					为什么是 1.41？因为 $\sqrt{2}\approx1.41$。%
				}。
				\item 使脚注编号和脚注文本之间默认间隔一个空格。
			\end{itemize}

			\item \textbf{\texttt{[no]tightlist}}：
				是否采用比 \LaTeX{} 默认设定更加紧密的枚举环境。
				在枚举环境（itemize、enumerate 和 description）中，
				每个条目的内容较少时，条目往往显得稀疏；
				在参考文献列表中也有类似的现象。
				启用 \verb|tightlist| 选项后，
				将去掉这些环境中额外增加的（垂直）间隔。

			\item \textbf{\texttt{[no]pdftoc}}\footnote{%
					此选项部分等价于 1.4 alpha2 及以前版本 pkuthss-extra 宏包%
					的 \texttt{[no]tocbibind} 选项。
					因为 tocbibind 宏包和 biblatex 宏包冲突，
					pkuthss-extra 宏包不再调用 tocbibind 宏包。%
				}：
				启用 \verb|pdftoc| 选项后，
				用 \verb|\tableofcontents| 命令生成目录时%
				会自动添加“目录”的 pdf 书签。

			\item \textbf{\texttt{[no]uppermark}}：
				是否在页眉中将章节名中的小写字母转换为大写字母。
				就目前而言，
				这样的转换存在着一些较为严重的缺陷\footnote{%
					准确地说是 \texttt{\string\MakeUppercase} 宏的问题，
					其在某些地方的转换不够健壮，
					例如 \texttt{\string\cite\string{ctex\string}}
					会被转换成 \texttt{\string\cite\string{CTEX\string}}。%
				}，
				因此不建议使用。
				基于上述考虑，%
				\myemph{%
					pkuthss-extra 宏包默认启用 \texttt{nouppermark} 选项，
					即在不在页眉中使用大写的章节名%
				}。

			\item \textbf{\texttt{[no]spechap}}\footnote{%
					“spechap”是“\textbf{spec}ial \textbf{chap}ter”的缩写。%
				}：
				是否启用第 \ref{ssec:misc} 小节中介绍的 %
				\verb|\specialchap| 命令。

			\item \textbf{\texttt{[no]pdfprop}}：
				是否自动根据设定的论文文档信息（如作者、标题等）
				设置生成的 pdf 文档的相应属性。%
				\myemph{%
					注意：
					该选项实际上是在 \texttt{\string\maketitle} 时生效的，
					这是因为考虑到%
					通常用户在调用 \texttt{\string\maketitle} 前%
					已经设置好所有的文档信息。
					若用户不调用 \texttt{\string\maketitle}，
					则需在设定完文档信息之后自行调用%
					第 \ref{ssec:misc} 小节中介绍的 %
					\texttt{\string\setpdfproperties} 命令以完成%
					pdf 文档属性的设定。%
				}

			\item \textbf{\texttt{[no]colorlinks}}\footnote{%
					此选项等价于 1.3 及以前版本 pkuthss-extra 宏包%
					的 \texttt{[no]linkcolor} 选项，
					但后来发现这会和 hyperref 宏包的一个同名选项冲突，
					故改为 \texttt{[no]colorlinks}。%
				}：
				是否在生成的 pdf 文档中使用彩色的链接。
		\end{itemize}

		例如，在提交打印版的论文时，
		彩色的链接文字在黑白打印出来之后可能颜色会很浅。
		此时用户\myemph{%
			可以启用 pkuthss-extra 宏包的 nocolorlinks 选项，
			使所有的链接变为黑色，以免影响打印%
		}：
\begin{Verbatim}[frame = single]
\documentclass[..., nocolorlinks]{pkuthss} % “...”代表其它的选项。
\end{Verbatim}
		用户还可以进一步修改 \verb|img/| 目录中 %
		\verb|pkulogo.eps| 和 \verb|pkuword.eps| 两个图片文件，
		以使封面上的北京大学图样也变为黑色（详见相应文件中的注释）。

	\section{pkuthss 文档模版提供的命令和环境}
		\subsection{设定文档信息的命令}

		这一类命令的语法为
\begin{Verbatim}[frame = single]
\commandname{具体信息} % commandname 为具体命令的名称。
\end{Verbatim}

		这些命令总结如下：
		\begin{itemize}
			\item \texttt{\bfseries\string\ctitle}：设定论文中文标题；
			\item \texttt{\bfseries\string\etitle}：设定论文英文标题；
			\item \texttt{\bfseries\string\cauthor}：设定作者的中文名；
			\item \texttt{\bfseries\string\eauthor}：设定作者的英文名；
			\item \texttt{\bfseries\string\studentid}：设定作者的学号；
			\item \texttt{\bfseries\string\date}：设定日期；
			\item \texttt{\bfseries\string\school}：设定作者的学院名；
			\item \texttt{\bfseries\string\cmajor}：设定作者专业的中文名；
			\item \texttt{\bfseries\string\emajor}：设定作者专业的英文名；
			\item \texttt{\bfseries\string\direction}：设定作者的研究方向；
			\item \texttt{\bfseries\string\cmentor}：设定导师的中文名；
			\item \texttt{\bfseries\string\ementor}：设定导师的英文名；
			\item \texttt{\bfseries\string\ckeywords}：设定中文关键词；
			\item \texttt{\bfseries\string\ekeywords}：设定英文关键词。
		\end{itemize}

		例如，如果要设定专业为“化学”（“Chemistry”），则可以使用以下命令：
\begin{Verbatim}[frame = single]
\cmajor{化学}
\emajor{Chemistry}
\end{Verbatim}

		\subsection{自身存储文档信息的命令}

		这一类命令的语法为
\begin{Verbatim}[frame = single]
% commandname 为具体的命令名。
\renewcommand{\commandname}{具体信息}
\end{Verbatim}

		这些命令总结如下：
		\begin{itemize}
			\item \texttt{\bfseries\string\cuniversity}：大学的中文名。
			\item \texttt{\bfseries\string\euniversity}：大学的英文名。
			\item \texttt{\bfseries\string\cthesisname}：论文类别的中文名。
			\item \texttt{\bfseries\string\ethesisname}：论文类别的英文名。
			\item \texttt{\bfseries\string\cabstractname}：摘要的中文标题。
			\item \texttt{\bfseries\string\eabstractname}：摘要的英文标题。
		\end{itemize}

		例如，
		如果要设定论文的类别为“本科生毕业论文”（“Undergraduate Thesis”），
		则可以使用以下命令：
\begin{Verbatim}[frame = single]
\renewcommand{\cthesisname}{本科生毕业论文}
\renewcommand{\ethesisname}{Undergraduate Thesis}
\end{Verbatim}

		\subsection{以“key = value”格式设置文档信息}

		用户可以通过 \verb|\pkuthssinfo| 命令集中设定文档信息，
		其语法为：
\begin{Verbatim}[frame = single]
% key1、key2、value1、value2 等为具体文档信息的项目名和内容。
\pkuthssinfo{key1 = value1, key2 = value2, ...}
\end{Verbatim}
		其中文档信息的项目名为前面提到的设定文档信息的命令名%
		或自身存储文档信息的命令名（不带反斜杠）。

		当文档信息的内容包含了逗号等有干扰的字符时，
		可以用大括号将这一项文档信息的全部内容括起来。%
		\myemph{%
			我们推荐用户总用大括号将文档信息的内容括起来，
			以避免很多不必要的麻烦。%
		}

		例如，前面提到的文档信息的设置可以集中地写成：
\begin{Verbatim}[frame = single, tabsize = 4]
\pkuthssinfo{
	..., % “...”代表其它的设定。
	cthesisname = {本科生毕业论文},
	ethesisname = {Undergraduate Thesis},
	cmajor = {化学}, emajor = {Chemistry}
}
\end{Verbatim}

		\subsection{pkuthss 文档模版提供的其它命令和环境\label{ssec:misc}}

		\texttt{\bfseries cabstract} 和 \texttt{\bfseries eabstract} %
		环境用于编写中英文摘要。
		用户只需要写摘要的正文；标题、作者、导师、专业等部分会自动生成。

		\texttt{\bfseries\string\specialchap} 命令%
		用于开始不进行标号但计入目录的一章，
		并合理安排其页眉。%
		\myemph{%
			注意：
			需要启用 pkuthss-extra 宏包的 \texttt{spechap} 选项%
			才能使用此命令。
			另外，在此章内的节或小节等命令应使用带星号的版本，
			例如 \texttt{\string\section\string*} 等，
			以免造成章节编号混乱。%
		}%
		例如，本文档中的“绪言”一章就是用 \verb|\specialchap{绪言}| %
		这条命令开始的。%

		\texttt{\bfseries\string\setpdfproperties} 命令%
		用于根据用户设定的文档信息自动设定生成的 pdf 文档的属性。
		此命令会在用户调用 \verb|\maketitle| 命令时被自动调用，
		因此通常不需要用户自己使用；
		但用户有时可能不需要输出标题页，
		从而不会调用 \verb|\maketitle| 命令，
		此时就需要在设定完文档信息之后调用 \verb|\setpdfproperties|。
		\myemph{%
			注意：
			需要启用 pkuthss-extra 宏包的 \texttt{pdfprop} 选项%
			才能使用此命令。%
		}

		\subsection{从其它文档类和宏包继承的功能}

		pkuthss 文档类建立在 ctexbook\supercite{ctex} 文档类的基础上，
		并调用了 CJKfntef、%
		graphicx\supercite{graphicx}、geometry\supercite{geometry}、%
		fancyhdr\supercite{fancyhdr} 和 %
		etoolbox\supercite{etoolbox} 等几个宏包。
		因此，ctexbook 文档类和这些宏包所提供的功能均可以使用。

		例如，用户如果想将目录的标题改为“目{\quad\quad}录”，
		则可以使用 ctexbook 文档类提供的 \verb|\CTEXoptions| 命令：
\begin{Verbatim}[frame = single]
\CTEXoptions{contentsname = {目{\quad\quad}录}}
\end{Verbatim}

		在默认的配置下，%
		pkuthss 文档模版使用作者编写的 %
		biblatex\supercite{biblatex} 样式\supercite{biblatex-caspervector}%
		进行参考文献和引用的排版，
		用户可以使用其提供的功能。
		例如，
		用户可以分别使用 \verb|\cite|、\verb|\parencite| 和 \verb|\supercite| %
		生成未格式化的、带方括号的和上标且带方括号的引用标记：
\begin{Verbatim}[frame = single]
\cite{ctex}，\parencite{ctex}，\supercite{ctex}
\end{Verbatim}
		在本文中将产生“\cite{ctex}，\parencite{ctex}，\supercite{ctex}”。

		pkuthss-extra 宏包可能调用以下这些宏包：
		\begin{itemize}
			\item 启用 \verb|spacing| 选项时会调用 %
				setspace 和 footmisc\supercite{footmisc} 宏包。
			\item 启用 \verb|tightlist| 选项时会调用 %
				enumitem\supercite{enumitem} 宏包。
		\end{itemize}
		因此在启用相应选项时，用户可以使用对应宏包所提供的功能。

		\subsection{不建议更改的设置}
		\myemph{%
			pkuthss 文档类中有一些一旦改动就有可能破坏预设排版规划的设置，
			因此不建议更改这些设置，它们是：
			\begin{itemize}
				\item 纸张类型：A4；
				\item 页芯尺寸：%
					$240\,\mathrm{mm}\times150\,\mathrm{mm}$，
					包含页眉、页脚；
				\item 默认字号：小四号。
			\end{itemize}%
		}

	\section{高级设置}\label{sec:advanced}

	pkuthss 文档模版的实现是简洁、清晰、灵活的。
	当一些细节的自定义无法通过模版提供的外部接口实现时，
	我们鼓励用户（在适当理解相关部分代码的前提下）通过修改模版进行自定义。

	一个常见的需求是封面中部分内容（特别是论文的标题、专业和研究方向）太长，
	超出了在预设的空间。
	此时，
	用户可以修改 \verb|pkuthss.cls| 里 \verb|\maketitle| 定义中
	\verb|\pkuthss@int@fillinblank| 宏的参数来改变
	带下划线的空白的行数和行宽，其语法为：
\begin{Verbatim}[frame = single]
\pkuthss@int@fillinblank{行数}{行宽}{内容}
\end{Verbatim}
	例如，如果“研究方向”一栏需要两行的空白，
	可以将 \verb|pkuthss.cls| 里的
\begin{Verbatim}[frame = single]
\pkuthss@int@fillinblank{1}{\pkuthss@tmp@len}{\kaishu\@direction}
\end{Verbatim}
	改为
\begin{Verbatim}[frame = single]
\pkuthss@int@fillinblank{2}{\pkuthss@tmp@len}{\kaishu\@direction}
\end{Verbatim}
	当然，为了美观，可以将多于一行的部分移到封面中作者信息部分的最下方。


	% vim:ts=4:sw=4
%
% Copyright (c) 2008-2009 solvethis
% Copyright (c) 2010-2012 Casper Ti. Vector
% Public domain.

\chapter{问题及其解决}
	\section{文档中已经提到的常见问题（按重要性排序）}

	在默认设置（启用 \verb|colorlinks| 选项）下，
	黑白打印时文档中的部分彩色链接可能会变成浅灰色，
	解决方式见第 \ref{ssec:extra} 小节。

	中文字体字库不全（只包含 GB2312 字符集内字符）时，
	生成的 pdf 文档中可能缺少部分字符，
	解决方式见第 \ref{sec:req} 节。

	\verb|img/| 目录中 eps 图片未转换为 pdf 格式时，%
	pdf\LaTeX{} 方式编译可能出错，
	解决方式见第 \ref{sec:doc-dir} 节。

	使用过旧的 \TeX{} 系统和各宏包，
	或使用某些 Linux 发行版软件仓库所提供的 \TeX{}Live 时，
	可能引起一些问题，
	详见第 \ref{sec:req} 节。

	文档默认情况下是双面模式，章末可能产生空白页，详见第 \ref{ssec:options} 小节。

	一些高级设置，
	如封面中部分内容长度超过预设空间容量时的设置，
	见第 \ref{sec:advanced} 节。

	biblatex 宏包\cite{biblatex}会自行设定 \verb|\bibname|，
	故会覆盖通过 \verb|\CTEXoptions| 设定的参考文献列表标题。
	使用 biblatex 的用户可以使用 \verb|\printbibliography| 的
	\verb|title| 选项来手动设定参考文献列表的标题，例如：
\begin{Verbatim}[frame = single]
\printbibliography[title = {文献}, ...] % “...”为其它选项。
\end{Verbatim}

	\section{其它可能存在的问题}
		\subsection{上游宏包可能引起的问题}

		hyperref 宏包\supercite{hyperref}和一些宏包可能发生冲突。
		关于如何避免这些冲突，可以参考 hyperref 宏包的文档。
		此文件通常和执行 \verb|texdoc hyperref| %
		时打开的 pdf 文件位于同一目录中。
		低于 1.02c 版本的 ctex 宏包中对 hyperref 的设置有些不周，
		因此文档类中对其进行了一些手动的处理。
		考虑到新版本 ctex 宏包将逐渐被更多人采用，
		进行这些处理的代码将在以后被删除，
		而改成直接调用 ctex 宏包的 \verb|hyperref| 选项。

		\subsection{文档格式可能存在的问题}

		研究生手册和其电子版\supercite{pku-thesisstyle}要求的论文封面并不一致。
		这里以电子版为准。

		\subsection{其它一些问题}

		使用 GBK 编码和 pdf\LaTeX{} 编译方式时需要用户%
		运行 \verb|gbk2uni| 程序来转换 \verb|.out| 文件，
		否则生成的 pdf 书签可能乱码。
		考虑到用户可能没有 \verb|gbk2uni| 程序，且有用户使用 UTF-8 编码，
		默认的 \verb|Makefile| 和 \verb|Make.bat| 中将相关代码注释掉了，
		用户可以自行去掉相应的注释。

	\section{反馈意见和建议}

	关于 pkuthss 文档模版的意见和建议，
	请在北大未名 BBS 的 MathTools 版或 %
	Google Code 上 pkuthss 项目的 issue tracker%
	\footnote{\url{http://code.google.com/p/caspervector/issues/list}.}%
	上提出，
	或通过电子邮件\footnote%
	{\href{mailto:CasperVector@gmail.com}{\texttt{CasperVector@gmail.com}}.}%
	告知 Casper Ti. Vector。
	上述三种反馈方法中，建议用户尽量采用靠前的方法。

	在进行反馈时，请尽量确保已经仔细阅读本文档中的说明。
	如果是通过 BBS 或电子邮件进行反馈，
	请在标题中说明是关于 pkuthss 文档模版的反馈；
	如果是通过 Google Code 进行反馈，
	请给 issue 加上 \verb|Proj-Pkuthss| 标签。
	如果是错误报告，
	请说明所使用 pkuthss 模版的版本、
	自己使用的操作系统和 \TeX{} 系统的类型和版本；
	同时强烈建议附上一个出错的最小例子及其相应的编译日志（\verb|.log| 文件），
	在文件较长时请使用附件。


	% 结论。
	% vim:ts=4:sw=4
%
% Copyright (c) 2008-2009 solvethis
% Copyright (c) 2010-2012 Casper Ti. Vector
% Public domain.

\specialchap{结论}

pkuthss 文档模版结构较为简洁、清晰、灵活，较为易于学习和使用。
希望它能为各位需要使用 \LaTeX{} 撰写学位论文的同学提供一些帮助。



	% 正文中的附录部分。
	\appendix
	% 排版参考文献列表，并使其出现在目录中。
	% 如果同时要使参考文献列表参与章节编号，可将“bibintoc”改为“bibnumbered”。
	\printbibliography[heading = bibintoc]
	% 各附录。
	% vim:ts=4:sw=4
%
% Copyright (c) 2008-2009 solvethis
% Copyright (c) 2010-2012 Casper Ti. Vector
% Public domain.

\chapter{pkuthss 文档模版的实现}
\raggedbottom % 避免某些奇怪的“Underfull \vbox”警告。

	\section{pkuthss 文档类和 pkuthss-extra 宏包的实现}	
		\subsection{共用文件头部}
		\VerbatimInput[
			frame = lines, fontsize = {\footnotesize}, tabsize = 2,
			baselinestretch = 1, lastline = 23, numbers = left
		]{pkuthss.cls}

		\subsection{\texttt{pkuthss.cls}}
		\VerbatimInput[
			frame = lines, fontsize = {\footnotesize}, tabsize = 2,
			baselinestretch = 1, firstline = 25, numbers = left
		]{pkuthss.cls}

		\subsection{\texttt{pkuthss-utf8.def} 和 \texttt{pkuthss-gbk.def}}
		\VerbatimInput[
			frame = lines, fontsize = {\footnotesize}, tabsize = 2,
			baselinestretch = 1, firstline = 25, numbers = left
		]{pkuthss-utf8.def}

		\subsection{\texttt{pkuthss-extra.sty}}
		\VerbatimInput[
			frame = lines, fontsize = {\footnotesize}, tabsize = 2,
			baselinestretch = 1, firstline = 25, numbers = left
		]{pkuthss-extra.sty}

	\section{pkuthss 说明（示例）文档的源代码}

	本文档的源代码中大部分已经有了较为详细的注释，
	故请直接参照相应文件中的注释。

	\myemph{%
		注：%
		\texttt{img/} 目录中的 \texttt{Makefile} 和%
		两个 PostScript（\texttt{.eps}）文件（都是文本文件）中%
		也有详细的注释哦 :)%
	}

\flushbottom % 取消 \raggedbottom 的作用。


	% vim:ts=4:sw=4
%
% Copyright (c) 2008-2009 solvethis
% Copyright (c) 2010-2012 Casper Ti. Vector
% Public domain.

\chapter{更新记录}
\raggedbottom % 避免某些奇怪的“Underfull \vbox”警告。

\section{1.3 版以后的更新记录}
\VerbatimInput[
	tabsize = 4, fontsize = {\small}, baselinestretch = 1
]{./verbatim/ChangeLog.txt}

\section{1.3 及其以前版本的更新记录}
\VerbatimInput[
	tabsize = 4, fontsize = {\small}, baselinestretch = 1.1
]{./verbatim/ChangeLog-upto-1.3.txt}

\flushbottom % 取消 \raggedbottom 的作用。



	% 以下为正文之后的部分。
	\backmatter

	% 致谢。
	% vim:ts=4:sw=4
%
% Copyright (c) 2008-2009 solvethis
% Copyright (c) 2010-2012 Casper Ti. Vector
% Public domain.

\chapter{致谢}

感谢北大未名 BBS 的 MathTools 版和 Thesis 版诸位同学的支持。
特别感谢 pkuthss 模版的最初创作者 solvethis 网友，
以及不断地对 Casper 提出的诸多问题予以解答的 cauchy 网友 :)


	% 原创性声明和使用授权说明。
	% vim:ts=4:sw=4
%
% Copyright (c) 2008-2009 solvethis
% Copyright (c) 2010-2012 Casper Ti. Vector
% All rights reserved.
%
% Redistribution and use in source and binary forms, with or without
% modification, are permitted provided that the following conditions are
% met:
%
% * Redistributions of source code must retain the above copyright notice,
%   this list of conditions and the following disclaimer.
% * Redistributions in binary form must reproduce the above copyright
%   notice, this list of conditions and the following disclaimer in the
%   documentation and/or other materials provided with the distribution.
% * Neither the name of Peking University nor the names of its contributors
%   may be used to endorse or promote products derived from this software
%   without specific prior written permission.
% 
% THIS SOFTWARE IS PROVIDED BY THE COPYRIGHT HOLDERS AND CONTRIBUTORS "AS
% IS" AND ANY EXPRESS OR IMPLIED WARRANTIES, INCLUDING, BUT NOT LIMITED TO,
% THE IMPLIED WARRANTIES OF MERCHANTABILITY AND FITNESS FOR A PARTICULAR
% PURPOSE ARE DISCLAIMED. IN NO EVENT SHALL THE COPYRIGHT HOLDER OR
% CONTRIBUTORS BE LIABLE FOR ANY DIRECT, INDIRECT, INCIDENTAL, SPECIAL,
% EXEMPLARY, OR CONSEQUENTIAL DAMAGES (INCLUDING, BUT NOT LIMITED TO,
% PROCUREMENT OF SUBSTITUTE GOODS OR SERVICES; LOSS OF USE, DATA, OR
% PROFITS; OR BUSINESS INTERRUPTION) HOWEVER CAUSED AND ON ANY THEORY OF
% LIABILITY, WHETHER IN CONTRACT, STRICT LIABILITY, OR TORT (INCLUDING
% NEGLIGENCE OR OTHERWISE) ARISING IN ANY WAY OUT OF THE USE OF THIS
% SOFTWARE, EVEN IF ADVISED OF THE POSSIBILITY OF SUCH DAMAGE.

% 原创性声明和使用授权说明页不需要装订到论文中，故不显示页码。
\cleardoublepage\thispagestyle{empty}
{
	\linespread{1.5}\selectfont
	\section*{原创性声明}

	本人郑重声明：
	所呈交的学位论文，是本人在导师的指导下，独立进行研究工作所取得的成果。
	除文中已经注明引用的内容外，
	本论文不含任何其他个人或集体已经发表或撰写过的作品或成果。
	对本文的研究做出重要贡献的个人和集体，均已在文中以明确方式标明。
	本声明的法律结果由本人承担。
	\vspace{2.5em}\par
	\rightline
	{%
		论文作者签名：\hspace{5em}%
		日期：\hspace{2em}年\hspace{2em}月\hspace{2em}日%
	}
    \par
}


	% 原创性声明和使用授权说明，不显示页码。
	%\pagestyle{empty}
	%% vim:ts=4:sw=4
%
% Copyright (c) 2008-2009 solvethis
% Copyright (c) 2010-2012 Casper Ti. Vector
% All rights reserved.
%
% Redistribution and use in source and binary forms, with or without
% modification, are permitted provided that the following conditions are
% met:
%
% * Redistributions of source code must retain the above copyright notice,
%   this list of conditions and the following disclaimer.
% * Redistributions in binary form must reproduce the above copyright
%   notice, this list of conditions and the following disclaimer in the
%   documentation and/or other materials provided with the distribution.
% * Neither the name of Peking University nor the names of its contributors
%   may be used to endorse or promote products derived from this software
%   without specific prior written permission.
% 
% THIS SOFTWARE IS PROVIDED BY THE COPYRIGHT HOLDERS AND CONTRIBUTORS "AS
% IS" AND ANY EXPRESS OR IMPLIED WARRANTIES, INCLUDING, BUT NOT LIMITED TO,
% THE IMPLIED WARRANTIES OF MERCHANTABILITY AND FITNESS FOR A PARTICULAR
% PURPOSE ARE DISCLAIMED. IN NO EVENT SHALL THE COPYRIGHT HOLDER OR
% CONTRIBUTORS BE LIABLE FOR ANY DIRECT, INDIRECT, INCIDENTAL, SPECIAL,
% EXEMPLARY, OR CONSEQUENTIAL DAMAGES (INCLUDING, BUT NOT LIMITED TO,
% PROCUREMENT OF SUBSTITUTE GOODS OR SERVICES; LOSS OF USE, DATA, OR
% PROFITS; OR BUSINESS INTERRUPTION) HOWEVER CAUSED AND ON ANY THEORY OF
% LIABILITY, WHETHER IN CONTRACT, STRICT LIABILITY, OR TORT (INCLUDING
% NEGLIGENCE OR OTHERWISE) ARISING IN ANY WAY OUT OF THE USE OF THIS
% SOFTWARE, EVEN IF ADVISED OF THE POSSIBILITY OF SUCH DAMAGE.

% 原创性声明和使用授权说明页不需要装订到论文中，故不显示页码。
\cleardoublepage\thispagestyle{empty}
{
	\linespread{1.5}\selectfont
	\section*{原创性声明}

	本人郑重声明：
	所呈交的学位论文，是本人在导师的指导下，独立进行研究工作所取得的成果。
	除文中已经注明引用的内容外，
	本论文不含任何其他个人或集体已经发表或撰写过的作品或成果。
	对本文的研究做出重要贡献的个人和集体，均已在文中以明确方式标明。
	本声明的法律结果由本人承担。
	\vspace{2.5em}\par
	\rightline
	{%
		论文作者签名：\hspace{5em}%
		日期：\hspace{2em}年\hspace{2em}月\hspace{2em}日%
	}
    \par
}


	%--成绩评定表--
	% \includepdf[pages=1-2]{chap/grade.pdf}
\end{document}

